\section{Two-Dimensional Wave and Heat Equations: Rectangle and Disc}
\def\EMfootchapter{Chapter 6 \ \textbullet\ Two-Dimensional Wave and Heat Equations: Rectangle and Disc}
\def\EMtitlesize{26}\def\EMtitleskip{31}
\begin{frame}{Chapter 6: Two-Dimensional Wave and Heat Equations: Rectangle and Disc}
{\fontsize{20}{26}\selectfont\bfseries\color{SLdef}\EMalign In this chapter\par}
\vspace{10bp}
\EMitem{\lr{\textbullet}}{\dimexpr\linewidth-26bp\relax}{The Two-Dimensional Wave Equation}
\EMitem{\lr{\textbullet}}{\dimexpr\linewidth-26bp\relax}{The Two-Dimensional Wave Equation with a Rectangular Boundary}
\EMitem{\lr{\textbullet}}{\dimexpr\linewidth-26bp\relax}{The Two-Dimensional Heat Equation with a Rectangular Boundary}
\EMitem{\lr{\textbullet}}{\dimexpr\linewidth-26bp\relax}{The Laplacian in Polar Coordinates}
\EMitem{\lr{\textbullet}}{\dimexpr\linewidth-26bp\relax}{The Two-Dimensional Wave Equation with a Circular Boundary}
\end{frame}

\def\EMtitlesize{26}\def\EMtitleskip{31}
\begin{frame}{The Two-Dimensional Wave Equation}
\begin{minipage}[t]{500bp}\vspace{0pt}
\begin{uncoverenv}<2->
\EMp{Modelling a vibrating membrane (deriving the equation governing \EMim{u(x,y,t)}).}
\end{uncoverenv}
\EMgap
\begin{uncoverenv}<3->
\EMp{\textbf{Physical assumptions:}}
\end{uncoverenv}
\EMgap
\begin{uncoverenv}<4->
\EMitem{\lr{\textbullet}}{\dimexpr\linewidth-26bp\relax}{The mass is distributed uniformly over the membrane (a homogeneous membrane), and the membrane is perfectly flexible.}
\end{uncoverenv}
\EMgap
\begin{uncoverenv}<5->
\EMitem{\lr{\textbullet}}{\dimexpr\linewidth-26bp\relax}{The membrane is stretched so tightly that the force of gravity on it is negligible.}
\end{uncoverenv}
\EMgap
\begin{uncoverenv}<6->
\EMitem{\lr{\textbullet}}{\dimexpr\linewidth-26bp\relax}{The deflection of the membrane from equilibrium is very small.}
\end{uncoverenv}
\EMgap
\end{minipage}
\EMside{230.6bp}{\uncover<7->{\EMfigure{1.500}{m06-membrane}{An element \EMim{\Delta x\times\Delta y} of the vibrating membrane and the tensile forces \EMim{T\Delta x} and \EMim{T\Delta y} acting on it}}}
\end{frame}

\def\EMtitlesize{26}\def\EMtitleskip{31}
\begin{frame}{The Two-Dimensional Wave Equation (continued)}
\begin{uncoverenv}<2->
\EMp{By a modelling process analogous to the one-dimensional case, the equation of the vibrating two-dimensional membrane is obtained:}
\vspace{3bp}
\EMeq{\frac{\partial^2u}{\partial t^2} =c^2\left(\frac{\partial^2u}{\partial x^2}+\frac{\partial^2u}{\partial y^2}\right) =c^2\nabla^2u}
\end{uncoverenv}
\EMgap
\begin{uncoverenv}<3->
\EMp{For a unique solution we need boundary conditions and initial conditions:}
\end{uncoverenv}
\EMgap
\begin{uncoverenv}<4->
\EMitem{\lr{\textbullet}}{\dimexpr\linewidth-26bp\relax}{Boundary conditions: \EMim{u(x,y,t)=0} on the boundary of the membrane}
\end{uncoverenv}
\EMgap
\begin{uncoverenv}<5->
\EMitem{\lr{\textbullet}}{\dimexpr\linewidth-26bp\relax}{Initial conditions: \EMim{u(x,y,0)=f(x,y)} and \EMim{u_t(x,y,0)=g(x,y)}, where \EMim{u_t\triangleq\dfrac{\partial u}{\partial t}}}
\end{uncoverenv}
\end{frame}

\def\EMtitlesize{26}\def\EMtitleskip{31}
\begin{frame}{The Two-Dimensional Wave Equation with a Rectangular Boundary}
\begin{minipage}[t]{500bp}\vspace{0pt}
\EMgap
\begin{uncoverenv}<3->
\EMeq{\frac{\partial^2u}{\partial t^2} =c^2\left(\frac{\partial^2u}{\partial x^2}+\frac{\partial^2u}{\partial y^2}\right)}
\end{uncoverenv}
\EMgap
\begin{uncoverenv}<4->
\EMitem{\lr{\textbullet}}{\dimexpr\linewidth-26bp\relax}{Boundary conditions: \EMim{u(0,y,t)=u(a,y,t)=0} and \EMim{u(x,0,t)=u(x,b,t)=0}}
\end{uncoverenv}
\EMgap
\begin{uncoverenv}<5->
\EMitem{\lr{\textbullet}}{\dimexpr\linewidth-26bp\relax}{Initial conditions: \EMim{u(x,y,0)=f(x,y)} and \EMim{u_t(x,y,0)=g(x,y)}}
\end{uncoverenv}
\end{minipage}
\EMside{0.0bp}{\uncover<2->{\EMfigure{1.500}{m06-rectangle}{The rectangular region \EMim{0\le x\le a}, \EMim{0\le y\le b}}}}
\end{frame}

\def\EMtitlesize{26}\def\EMtitleskip{31}
\begin{frame}{Solving the Equation}
\begin{uncoverenv}<2->
\EMp{Here we solve the equation by \textbf{separation of variables}. Assume \EMim{u(x,y,t)=X(x)\times Y(y)\times T(t)}.}
\vspace{3bp}
\EMeq{\frac{\partial^2u}{\partial x^2} =X''YT,\qquad \frac{\partial^2u}{\partial y^2} =XY''T,\qquad \frac{\partial^2u}{\partial t^2} =XYT''\quad \Longrightarrow\quad XYT'' =c^2\big(X''YT+XY''T\big)}
\end{uncoverenv}
\EMgap
\begin{uncoverenv}<3->
\EMeq{\frac{T''}{c^2T} =\frac{X''}X+\frac{Y''}Y =-\lambda^2\  \Rightarrow\ \frac{Y''}Y =-\lambda^2-\frac{X''}X =-\rho^2\  \Rightarrow\ \frac{X''}X =-\lambda^2+\rho^2 =-\mu^2}
\end{uncoverenv}
\EMgap
\begin{uncoverenv}<4->
\EMeq{\begin{cases}X''+\mu^2X=0\\ Y''+\rho^2Y=0\\ T''+c^2\lambda^2T=0\end{cases}\qquad \lambda^2 =\mu^2+\rho^2}
\end{uncoverenv}
\EMgap
\begin{uncoverenv}<5->
\EMp{\textbf{The solution in the \EMim{x} direction:}}
\vspace{3bp}
\EMeq{X(x) =A_1\cos\mu x+A_2\sin\mu x,\quad X(0) =0 \Rightarrow A_1 =0,\quad X(a) =0 \Rightarrow A_2\sin\mu a =0 \Rightarrow\mu a =n\pi \Rightarrow\mu_n \\[4pt] =\frac{n\pi}a}
\end{uncoverenv}
\EMgap
\begin{uncoverenv}<6->
\EMeq{X_n(x) =A_n\sin\mu_nx,\qquad \mu_n =\frac{n\pi}a}
\end{uncoverenv}
\end{frame}

\def\EMtitlesize{26}\def\EMtitleskip{31}
\begin{frame}{Solving the Equation (continued)}
\begin{uncoverenv}<2->
\EMp{\textbf{The solution in the \EMim{y} direction:}}
\vspace{3bp}
\EMeq{Y(y) =B_1\cos\rho y+B_2\sin\rho y,\quad Y(0) =0 \Rightarrow B_1 =0,\quad Y(b) =0 \Rightarrow B_2\sin\rho b =0 \Rightarrow\rho b =m\pi \Rightarrow\rho_m \\[4pt] =\frac{m\pi}b}
\end{uncoverenv}
\EMgap
\begin{uncoverenv}<3->
\EMeq{Y_m(y) =B_m\sin\rho_my,\qquad \rho_m =\frac{m\pi}b}
\end{uncoverenv}
\EMgap
\begin{uncoverenv}<4->
\EMp{\textbf{The solution in time:}}
\vspace{3bp}
\EMeq{\lambda_{nm}^2 =\mu_n^2+\rho_m^2 =\left(\frac{n\pi}a\right)^2+\left(\frac{m\pi}b\right)^2,\qquad T_{nm}(t) =C_1\cos(c\lambda_{nm}t)+C_2\sin(c\lambda_{nm}t)}
\end{uncoverenv}
\EMgap
\begin{uncoverenv}<5->
\EMeq{u_{nm}(x,y,t) =\big[A_{nm}\cos(c\lambda_{nm}t)+B_{nm}\sin(c\lambda_{nm}t)\big]\sin\!\left(\frac{n\pi}ax\right)\sin\!\left(\frac{m\pi}by\right),\qquad n,m =1,2,\dots}
\end{uncoverenv}
\EMgap
\begin{uncoverenv}<6->
\EMeq{u(x,y,t) =\sum_{n=1}^{\infty}\sum_{m=1}^{\infty}\big[A_{nm}\cos(c\lambda_{nm}t)+B_{nm}\sin(c\lambda_{nm}t)\big]\sin\!\left(\frac{n\pi}ax\right)\sin\!\left(\frac{m\pi}by\right)}
\end{uncoverenv}
\end{frame}

\def\EMtitlesize{26}\def\EMtitleskip{31}
\begin{frame}{Solving the Equation (continued)}
\begin{uncoverenv}<2->
\EMp{\textbf{Applying the initial condition \EMim{u(x,y,0)=f(x,y)}:}}
\vspace{3bp}
\EMeq{f(x,y) =\sum_{n=1}^{\infty}\sum_{m=1}^{\infty}A_{nm}\sin\!\left(\frac{n\pi}ax\right)\sin\!\left(\frac{m\pi}by\right) =\sum_{n=1}^{\infty}K_n(y)\sin\frac{n\pi}ax,\qquad K_n(y)\triangleq\sum_{m=1}^{\infty}A_{nm}\sin\frac{m\pi}by}
\end{uncoverenv}
\EMgap
\begin{uncoverenv}<3->
\EMeq{K_n(y) =\frac2a\int_0^af(x,y)\sin\frac{n\pi}ax\,dx,\qquad A_{nm} =\frac2b\int_0^bK_n(y)\sin\frac{m\pi}by\,dy}
\end{uncoverenv}
\EMgap
\begin{uncoverenv}<4->
\EMeq{A_{nm} =\frac4{ab}\int_0^b\left\{\int_0^af(x,y)\sin\frac{n\pi}ax\,dx\right\}\sin\frac{m\pi}by\,dy}
\end{uncoverenv}
\end{frame}

\def\EMtitlesize{26}\def\EMtitleskip{31}
\begin{frame}{Solving the Equation (continued)}
\begin{EMbox}{exe}{1}{}
\EMp{Prove that}
\vspace{3bp}
\EMeq{B_{nm} =\frac4{abc\lambda_{nm}}\int_0^b\left\{\int_0^ag(x,y)\sin\frac{n\pi}ax\,dx\right\}\sin\frac{m\pi}by\,dy}
\end{EMbox}
\begin{EMbox}{imp}{1}{}
\begin{uncoverenv}<2->
\EMeq{u(x,y,t) =\sum_{n=1}^{\infty}\sum_{m=1}^{\infty}\big[A_{nm}\cos(c\lambda_{nm}t)+B_{nm}\sin(c\lambda_{nm}t)\big]\sin\!\left(\frac{n\pi}ax\right)\sin\!\left(\frac{m\pi}by\right),\qquad \lambda_{nm}^2 \\[4pt] =\left(\frac{n\pi}a\right)^2+\left(\frac{m\pi}b\right)^2}
\end{uncoverenv}
\EMgap
\begin{uncoverenv}<3->
\EMeq{A_{nm} =\frac4{ab}\int_0^b\left\{\int_0^af(x,y)\sin\frac{n\pi}ax\,dx\right\}\sin\frac{m\pi}by\,dy,\qquad B_{nm} \\[4pt] =\frac4{abc\lambda_{nm}}\int_0^b\left\{\int_0^ag(x,y)\sin\frac{n\pi}ax\,dx\right\}\sin\frac{m\pi}by\,dy}
\end{uncoverenv}
\end{EMbox}
\end{frame}

\def\EMtitlesize{26}\def\EMtitleskip{31}
\begin{frame}{Solving the Equation (continued)}
\begin{EMbox}{ex}{1}{}
\EMp{Find the solution of the two-dimensional motion equation above for \EMim{g(x,y)=0} and \EMim{f(x,y)=xy}.}
\EMgap
\begin{uncoverenv}<2->
\EMeq{g(x,y) =0\  \Rightarrow\ B_{nm} =0}
\end{uncoverenv}
\EMgap
\begin{uncoverenv}<3->
\EMeq{A_{nm} =\frac4{ab}\int_0^b\left\{\int_0^axy\sin\frac{n\pi}ax\,dx\right\}\sin\frac{m\pi}by\,dy =\frac4{ab}\left\{\left[\int_0^by\sin\frac{m\pi}by\,dy\right]\times\left[\int_0^ax\sin\frac{n\pi}ax\,dx\right]\right\}}
\end{uncoverenv}
\EMgap
\begin{uncoverenv}<4->
\EMeq{=\frac4{ab}\left[y\left(\frac{-b}{m\pi}\cos\frac{m\pi}by\right)-(1)\left(\frac{-b^2}{m^2\pi^2}\sin\frac{m\pi}by\right)\right]_0^b\times\left[x\left(\frac{-a}{n\pi}\cos\frac{n\pi}ax\right)-(1)\left(\frac{-a^2}{n^2\pi^2}\sin\frac{n\pi}ax\right)\right]_0^a}
\end{uncoverenv}
\EMgap
\begin{uncoverenv}<5->
\EMeq{=\frac4{ab}\cdot\frac{-b^2}{m\pi}(-1)^m\cdot\frac{-a^2}{n\pi}(-1)^n =\frac{4ab}{nm\pi^2}(-1)^{n+m}}
\end{uncoverenv}
\EMgap
\begin{uncoverenv}<6->
\EMeq{u(x,y,t) =\sum_{n=1}^{\infty}\sum_{m=1}^{\infty}\frac{4ab}{nm\pi^2}(-1)^{n+m}\cos(c\lambda_{nm}t)\sin\!\left(\frac{n\pi}ax\right)\sin\!\left(\frac{m\pi}by\right)}
\end{uncoverenv}
\end{EMbox}
\end{frame}

\def\EMtitlesize{26}\def\EMtitleskip{31}
\begin{frame}{The Two-Dimensional Heat Equation with a Rectangular Boundary}
\begin{uncoverenv}<2->
\EMeq{\frac{\partial u}{\partial t} =c^2\left(\frac{\partial^2u}{\partial x^2}+\frac{\partial^2u}{\partial y^2}\right) =c^2\nabla^2u}
\end{uncoverenv}
\EMgap
\begin{uncoverenv}<3->
\EMitem{\lr{\textbullet}}{\dimexpr\linewidth-26bp\relax}{Boundary condition: \EMim{u(x,y,t)=0} on the boundary of the plate}
\end{uncoverenv}
\EMgap
\begin{uncoverenv}<4->
\EMitem{\lr{\textbullet}}{\dimexpr\linewidth-26bp\relax}{Initial condition: \EMim{u(x,y,0)=f(x,y)}}
\end{uncoverenv}
\end{frame}

\def\EMtitlesize{26}\def\EMtitleskip{31}
\begin{frame}{Solving the Equation}
\begin{uncoverenv}<2->
\EMp{Assume \EMim{u(x,y,t)=X(x)\times Y(y)\times T(t)}.}
\vspace{3bp}
\EMeq{\frac{\partial^2u}{\partial x^2} =X''YT,\qquad \frac{\partial^2u}{\partial y^2} =XY''T,\qquad \frac{\partial u}{\partial t} =XYT'\quad \Longrightarrow\quad XYT' =c^2\big(X''YT+XY''T\big)}
\end{uncoverenv}
\EMgap
\begin{uncoverenv}<3->
\EMeq{\frac{T'}{c^2T} =\frac{X''}X+\frac{Y''}Y =-\lambda^2,\qquad \frac{Y''}Y =-\lambda^2-\frac{X''}X =-\rho^2,\qquad \frac{X''}X =-\lambda^2+\rho^2 =-\mu^2}
\end{uncoverenv}
\EMgap
\begin{uncoverenv}<4->
\EMeq{\begin{cases}X''+\mu^2X=0\\ Y''+\rho^2Y=0\\ T'+c^2\lambda^2T=0\end{cases}\qquad \lambda^2 =\mu^2+\rho^2}
\end{uncoverenv}
\EMgap
\begin{uncoverenv}<5->
\EMp{As in the wave case (\EMim{X_n(x)=A_n\sin\mu_nx} with \EMim{\mu_n=\frac{n\pi}a} and \EMim{Y_m(y)=B_m\sin\rho_my} with \EMim{\rho_m=\frac{m\pi}b}), and with}
\vspace{3bp}
\EMeq{T_{nm}(t) =C_1e^{-c^2\lambda_{nm}^2t},\qquad \lambda_{nm}^2 =\mu_n^2+\rho_m^2 =\left(\frac{n\pi}a\right)^2+\left(\frac{m\pi}b\right)^2}
\end{uncoverenv}
\EMgap
\begin{uncoverenv}<6->
\EMeq{u_{nm}(x,y,t) =A_{nm}\sin\!\left(\frac{n\pi}ax\right)\sin\!\left(\frac{m\pi}by\right)e^{-c^2\lambda_{nm}^2t},\qquad n =1,2,\dots,\ m =1,2,\dots}
\end{uncoverenv}
\end{frame}

\def\EMtitlesize{26}\def\EMtitleskip{31}
\begin{frame}{Solving the Equation (continued)}
\begin{uncoverenv}<2->
\EMeq{u(x,y,t) =\sum_{n=1}^{\infty}\sum_{m=1}^{\infty}A_{nm}\sin\!\left(\frac{n\pi}ax\right)\sin\!\left(\frac{m\pi}by\right)e^{-c^2\lambda_{nm}^2t}}
\end{uncoverenv}
\EMgap
\begin{uncoverenv}<3->
\EMp{\textbf{Applying the initial condition:} exactly as in the wave case,}
\vspace{3bp}
\EMeq{A_{nm} =\frac4{ab}\int_0^b\left\{\int_0^af(x,y)\sin\frac{n\pi}ax\,dx\right\}\sin\frac{m\pi}by\,dy}
\end{uncoverenv}
\end{frame}

\def\EMtitlesize{26}\def\EMtitleskip{31}
\begin{frame}{The Laplacian in Polar Coordinates}
\begin{minipage}[t]{500bp}\vspace{0pt}
\begin{uncoverenv}<2->
\EMp{The aim is to write \EMim{\nabla^2u=u_{xx}+u_{yy}} in polar coordinates in terms of \EMim{r,\theta}.}
\end{uncoverenv}
\EMgap
\EMgap
\begin{uncoverenv}<4->
\EMeq{r =\sqrt{x^2+y^2},\quad x =r\cos\theta,\qquad \theta =\tan^{-1}\!\left(\frac yx\right),\quad y \\[4pt] =r\sin\theta}
\end{uncoverenv}
\EMgap
\begin{uncoverenv}<5->
\EMeq{u_x =u_rr_x+u_\theta\theta_x\  \Rightarrow\ u_{xx} =\big(u_rr_x+u_\theta\theta_x\big)_x \\[4pt] =u_{rr}r_x^2+u_{\theta\theta}\theta_x^2+2r_x\theta_xu_{r\theta}+r_{xx}u_r+\theta_{xx}u_\theta}
\end{uncoverenv}
\EMgap
\begin{uncoverenv}<6->
\EMeq{u_y =u_rr_y+u_\theta\theta_y\  \Rightarrow\ u_{yy} \\[4pt] =u_{rr}r_y^2+u_{\theta\theta}\theta_y^2+2r_y\theta_yu_{r\theta}+r_{yy}u_r+\theta_{yy}u_\theta}
\end{uncoverenv}
\EMgap
\begin{uncoverenv}<7->
\EMeq{r_x =\frac xr,\quad r_{xx} =\frac{y^2}{r^3},\quad \theta_x =\frac{-y}{r^2},\quad \theta_{xx} =\frac{2xy}{r^4},\qquad r_y \\[4pt] =\frac yr,\quad r_{yy} =\frac{x^2}{r^3},\quad \theta_y =\frac x{r^2},\quad \theta_{yy} =\frac{-2xy}{r^4}}
\end{uncoverenv}
\end{minipage}
\EMside{57.6bp}{\uncover<3->{\EMfigure{1.500}{m06-polar}{Polar coordinates: \EMim{x=r\cos\theta}, \EMim{y=r\sin\theta}}}}
\end{frame}

\def\EMtitlesize{26}\def\EMtitleskip{31}
\begin{frame}{The Laplacian in Polar Coordinates (continued)}
\begin{uncoverenv}<2->
\EMeq{u_{xx} =\frac{x^2}{r^2}u_{rr}+\frac{y^2}{r^4}u_{\theta\theta}+\frac{-2xy}{r^3}u_{r\theta}+\frac{y^2}{r^3}u_r+\frac{2xy}{r^4}u_\theta}
\end{uncoverenv}
\EMgap
\begin{uncoverenv}<3->
\EMeq{u_{yy} =\frac{y^2}{r^2}u_{rr}+\frac{x^2}{r^4}u_{\theta\theta}+\frac{2xy}{r^3}u_{r\theta}+\frac{x^2}{r^3}u_r+\frac{-2xy}{r^4}u_\theta}
\end{uncoverenv}
\EMgap
\begin{uncoverenv}<4->
\EMp{Adding these two relations:}
\end{uncoverenv}
\EMgap
\begin{EMbox}{imp}{4}{}
\begin{uncoverenv}<5->
\EMeq{\nabla^2u =u_{rr}+\frac1{r^2}u_{\theta\theta}+\frac1ru_r}
\end{uncoverenv}
\end{EMbox}
\end{frame}

\def\EMtitlesize{26}\def\EMtitleskip{31}
\begin{frame}{The Two-Dimensional Wave Equation with a Circular Boundary}
\begin{minipage}[t]{500bp}\vspace{0pt}
\EMgap
\begin{uncoverenv}<3->
\EMeq{\frac{\partial^2u}{\partial t^2} =c^2\nabla^2u,\qquad u(R,t) =0\ \ \text{(boundary condition)},\qquad \\[4pt] u(r,0) =f(r),\ \ u_t(r,0) =g(r)\ \ \text{(initial conditions)}}
\end{uncoverenv}
\end{minipage}
\EMside{0.0bp}{\uncover<2->{\EMfigure{1.500}{m06-disc}{A circular membrane of radius \EMim{R}}}}
\end{frame}

\def\EMtitlesize{26}\def\EMtitleskip{31}
\begin{frame}{Solving the Equation}
\begin{uncoverenv}<2->
\EMp{Because of the circular symmetry of this problem it is better to solve it in polar coordinates. On a circular membrane the vibration is not a function of \EMim{\theta}, so}
\vspace{3bp}
\EMeq{\begin{cases}\dfrac{\partial^2u}{\partial t^2}=c^2\left(u_{rr}+\dfrac1ru_r\right)\\\noalign{\vskip 2mm} u(R,t)=0,\quad u(r,0)=f(r),\quad u_t(r,0)=g(r)\end{cases}}
\end{uncoverenv}
\EMgap
\begin{uncoverenv}<3->
\EMp{Now, assuming separable variables \EMim{u(r,t)=P(r)\times T(t)}:}
\vspace{3bp}
\EMeq{PT'' =c^2\left(P''T+\frac1rP'T\right)\  \Rightarrow\ \frac{T''}{c^2T} =\frac{P''}P+\frac1r\frac{P'}P =-\lambda^2\  \Rightarrow\ \begin{cases}P''+\dfrac1rP'+\lambda^2P=0\\\noalign{\vskip 1mm} T''+c^2\lambda^2T=0\end{cases}}
\end{uncoverenv}
\end{frame}

\def\EMtitlesize{26}\def\EMtitleskip{31}
\begin{frame}{Bessel's Differential Equation of Order \EMim{n}}
\begin{minipage}[t]{500bp}\vspace{0pt}
\begin{uncoverenv}<2->
\EMeq{y''+\frac1xy'+\left(1-\frac{n^2}{x^2}\right)y =0\quad \Longrightarrow\quad y(x) \\[4pt] =AJ_n(x)+BY_n(x)}
\end{uncoverenv}
\EMgap
\begin{uncoverenv}<3->
\EMp{where \EMim{J_n} are the Bessel functions of the first kind of order \EMim{n} and \EMim{Y_n} the Bessel functions of the second kind of order \EMim{n}.}
\end{uncoverenv}
\EMgap
\end{minipage}
\EMside{156.4bp}{\uncover<4->{\EMfigure{0.880}{m06-bessel}{The Bessel functions of the first kind \EMim{J_0,\dots,J_4} and of the second kind \EMim{Y_0,\dots,Y_4} for \EMim{0<x\le10}}}}
\end{frame}

\def\EMtitlesize{26}\def\EMtitleskip{31}
\begin{frame}{Bessel's Differential Equation of Order \EMim{n} (continued)}
\begin{uncoverenv}<2->
\EMp{With the substitution \EMim{s=\lambda r}:}
\vspace{3bp}
\EMeq{P' =\frac{dP}{dr} =\frac{dP}{ds}\frac{ds}{dr} =\lambda\frac{dP}{ds},\qquad P'' =\lambda^2\frac{d^2P}{ds^2},\qquad \frac1r\lambda =\frac{\lambda^2}s}
\end{uncoverenv}
\EMgap
\begin{uncoverenv}<3->
\EMeq{\lambda^2\frac{d^2P}{ds^2}+\frac1r\lambda\frac{dP}{ds}+\lambda^2P =0\  \Rightarrow\ \frac{d^2P}{ds^2}+\frac1s\frac{dP}{ds}+P =0}
\end{uncoverenv}
\EMgap
\begin{uncoverenv}<4->
\EMp{This is Bessel's equation of order zero (\EMim{n=0}), so}
\vspace{3bp}
\EMeq{P(r) =CJ_0(\lambda r)+DY_0(\lambda r)}
\end{uncoverenv}
\EMgap
\begin{uncoverenv}<5->
\EMp{\textbf{Applying the boundary conditions:}}
\end{uncoverenv}
\EMgap
\begin{uncoverenv}<6->
\EMitem{\lr{\textbullet}}{\dimexpr\linewidth-26bp\relax}{Since \EMim{P(0)\neq\infty} and \EMim{Y_0} is unbounded at the origin, the coefficient of \EMim{Y_0} is zero (\EMim{D=0}) and \EMim{P(r)=CJ_0(\lambda r)}.}
\end{uncoverenv}
\end{frame}

\def\EMtitlesize{26}\def\EMtitleskip{31}
\begin{frame}{Bessel's Differential Equation of Order \EMim{n} (continued)}
\begin{minipage}[t]{500bp}\vspace{0pt}
\begin{uncoverenv}<2->
\EMitem{\lr{\textbullet}}{\dimexpr\linewidth-26bp\relax}{\EMim{P(R)=0} gives \EMim{J_0(\lambda R)=0}, i.e.~\EMim{\lambda R=\alpha_m}, where \EMim{\alpha_m} are the roots of \EMim{J_0}:}
\vspace{3bp}
\EMeq{\lambda =\lambda_m =\frac{\alpha_m}R,\quad m =1,2,\dots\quad \Longrightarrow\quad P_m(r) \\[4pt] =C_mJ_0\!\left(\frac{\alpha_m}Rr\right)}
\end{uncoverenv}
\EMgap
\EMgap
\begin{uncoverenv}<4->
\EMeq{T''+c^2\lambda_m^2T =0\  \Rightarrow\ T_m(t) =A_m\cos(c\lambda_mt)+B_m\sin(c\lambda_mt)}
\end{uncoverenv}
\EMgap
\begin{EMbox}{imp}{4}{}
\begin{uncoverenv}<5->
\EMeq{u(r,t) \\[4pt] =\sum_{m=1}^{\infty}\big(A_m\cos(c\lambda_mt)+B_m\sin(c\lambda_mt)\big)J_0\!\left(\frac{\alpha_m}Rr\right),\qquad \\[4pt] \lambda_m =\frac{\alpha_m}R}
\end{uncoverenv}
\end{EMbox}
\end{minipage}
\EMside{154.9bp}{\uncover<3->{\EMfigure{1.283}{m06-j0-zeros}{The function \EMim{J_0(x)} and its roots \EMim{\alpha_1,\alpha_2,\alpha_3}}}}
\end{frame}

\def\EMtitlesize{26}\def\EMtitleskip{31}
\begin{frame}{Bessel's Differential Equation of Order \EMim{n} (continued)}
\begin{EMbox}{rem}{1}{}
\begin{uncoverenv}<2->
\EMp{Prepare a complete report on Bessel's differential equation of order \EMim{n}, the Bessel functions of the first and second kind and their properties, as well as the complete solution of the two-dimensional wave equation with a circular boundary (the report must be typed, in Word format).}
\end{uncoverenv}
\end{EMbox}
\end{frame}
